% BHU PYQ -- Auto-generated & error-corrected
\documentclass[11pt,a4paper]{article}

\usepackage[a4paper,top=2.5cm,bottom=2.5cm,left=2.8cm,right=2.8cm,headheight=14pt]{geometry}
\usepackage{amsmath,amssymb}
\usepackage[T1]{fontenc}
\usepackage[utf8]{inputenc}
\usepackage{lmodern}
\usepackage{microtype}
\usepackage{fancyhdr}
\usepackage{enumitem}
\usepackage{hyperref}
\usepackage{lastpage}
\usepackage{parskip}

\hypersetup{colorlinks=true,urlcolor=black,linkcolor=black,
  pdftitle={BPT-201: Thermal Physics -- BHU PYQ},pdfauthor={Banaras Hindu University}}

\pagestyle{fancy}\fancyhf{}
\renewcommand{\headrulewidth}{0.4pt}\renewcommand{\footrulewidth}{0.4pt}
\fancyhead[L]{\small   \textit{Banaras Hindu University}}
\fancyhead[R]{\small   \textit{BPT-201: Thermal Physics}}
\fancyfoot[C]{\small Page  \thepage\ of \pageref{LastPage}}


\newlist{parts}{enumerate}{2}
\setlist[parts,1]{label=(\alph*),leftmargin=*,itemsep=5pt,topsep=3pt}
\setlist[parts,2]{label=(\roman*),leftmargin=*,itemsep=3pt,topsep=2pt}

\newcommand{\pts}[1]{\hfill   \textit{[#1]}}


\begin{document}


\begin{center}
  {\large\bfseries BANARAS HINDU UNIVERSITY}\\[2pt]
  {
\normalsize B.Sc. (Hons.) Semester II Examination, 2023-24}\\[6pt]
  \rule{0.8\linewidth}{0.5pt}\\[6pt]
  {\large\bfseries Paper No. BPT-201: Thermal Physics}\\[4pt]
  {
\normalsize Time: 3 Hours\hfill Maximum Marks: 70}
\end{center}
\rule{\linewidth}{0.4pt}
\smallskip



\noindent   \textit{Instructions: Answer   \textbf{five} questions including Question~1 which is compulsory. Symbols have their usual meanings.}
\bigskip

\medskip


\noindent   \textbf{Question 1.} Answer any   \textbf{five} of the following:

\begin{parts}
  \item What do you mean by inversion temperature? How does it depend on van der Waals gas constants? \pts{4}
  \item Prove the following relation for a Carnot cycle:
  \[
  \oint P \, dV = \oint T \, dS
  \] \pts{4}
  \item Prove that for an isothermal and isobaric process, $dG = 0$, where $G$ is Gibbs' free energy. \pts{4}
  \item Explain the role of temperature in superconductivity. Also explain the behaviour of a superconductor in a magnetic field (Meissner effect). \pts{4}
  \item Explain Kelvin's thermodynamical scale of temperature and define the absolute zero temperature. \pts{4}
  \item Define the efficiency of a Carnot's engine. Find its value if the engine is working between the temperatures of steam ($100^\circ \text{C}$) and ice ($0^\circ \text{C}$). If it absorbs 80 cal of heat, find the amount of heat rejected. \pts{4}
  \item Write down Wien's displacement law and find the wavelength at which human beings radiate maximum energy. (Wien's constant is $0.3 \text{ cm-K}$, human body temperature $\approx 37^\circ \text{C}$). \pts{4}
\end{parts}
\medskip\rule{\linewidth}{0.2pt}

\medskip


\noindent   \textbf{Question 2.}

\begin{parts}
  \item What do you mean by thermodynamic potentials? By using these potentials, obtain the following Maxwell thermodynamic relations:
  \[
  \left(\frac{\partial T}{\partial V}\right)_S = -\left(\frac{\partial P}{\partial S}\right)_V \quad  \text{and} \quad \left(\frac{\partial T}{\partial P}\right)_S = \left(\frac{\partial V}{\partial S}\right)_P
  \] \pts{8}
  \item Calculate the change in the boiling point of water when the pressure is increased by 1 atmosphere. (Normal boiling temperature of water = $100^\circ \text{C}$; Specific volume of steam = $1.671 \text{ m}^3 \text{kg}^{-1}$; Latent heat of steam = $2.268  \times 10^6 \text{ J kg}^{-1}$; 1 atm = $1.013  \times 10^5 \text{ N/m}^2$). \pts{6}
\end{parts}
\medskip\rule{\linewidth}{0.2pt}

\medskip


\noindent   \textbf{Question 3.}

\begin{parts}
  \item Explain the difference between an ideal and a real gas. Deduce the expression for the Joule-Thomson coefficient and explain the cooling and heating effects for a real gas. \pts{10}
  \item Calculate the critical constants $P_c$, $V_c$, and $T_c$ in terms of van der Waals constants $a$ and $b$. \pts{4}
\end{parts}
\medskip\rule{\linewidth}{0.2pt}

\hfill   \textbf{OR}


\begin{parts}
  \item Describe Kamerlingh Onnes' method of liquefaction of Helium gas with proper details of the apparatus used. Explain the changes in the various properties of liquid helium at temperatures below $2.19 \text{ K}$ ($\lambda$-transition). \pts{8}
  \item Describe the principle of regenerative cooling with a schematic diagram of the setup. \pts{6}
\end{parts}
\medskip\rule{\linewidth}{0.2pt}

\medskip


\noindent   \textbf{Question 4.}

\begin{parts}
  \item Write the Maxwell's thermodynamic relations and by using them derive the following relation:
  \[
  C_p - C_v = -T \left(\frac{\partial V}{\partial T}\right)_P^2 \left(\frac{\partial P}{\partial V}\right)_T
  \] \pts{8}
  \item Find the mean free path and the frequency of collision of gas molecules if the coefficient of viscosity of the gas at NTP is $1.66  \times 10^{-4} \text{ poises}$. (Use the gas parameters $v_{rms} = 450 \text{ m/s}$, and density $\rho = 1.25 \text{ g/lit}$). \pts{6}
\end{parts}
\medskip\rule{\linewidth}{0.2pt}

\hfill   \textbf{OR}


\begin{parts}
  \item Derive the Clausius-Clapeyron equation for a first-order phase transition and explain the response of Gibbs' function $G$ and its first derivatives w.r.t. temperature and pressure at first-order phase transition. \pts{8}
  \item Explain the ultraviolet catastrophe in black-body radiation. \pts{6}
\end{parts}
\medskip\rule{\linewidth}{0.2pt}

\medskip


\noindent   \textbf{Question 5.}

\begin{parts}
  \item Derive the following latent heat equation for saturated vapour:
  \[
  C_{sat,2} - C_{sat,1} = \frac{dL}{dT} - \frac{L}{T}
  \]
  where $C_{sat,2}$ and $C_{sat,1}$ are the specific heats of saturated vapour and saturated liquid respectively. \pts{7}
  \item Derive the following relation for the entropy of a perfect gas:
  \[
  S = C_v \log_e T + (C_p - C_v) \log_e V +  \text{constant}
  \] \pts{7}
\end{parts}

\vfill
\rule{\linewidth}{0.4pt}

\begin{center}\small
  \href{https://bhu-exam.vercel.app/}{Click here to visit our website}\\[4pt]
  \href{https://me-aryan.vercel.app/}{Click here to visit our contributor}\\[4pt]
  \href{https://quashan-me.vercel.app/}{Click here to visit our contributor}
\end{center}

\end{document}